
structure de \(\lambda\)-anneau sur un anneau \(A\), il revient au même de se donner un homomorphisme d'anneaux \(A\to\Ghat(A)_o\), tel que le diagramme 3.7.1 soit commutatif, ou un homomorphisme d'anneaux \(A\to\Ghat(A)_\nabla\), tel que le diagramme 3.7.2 soit commutatif~; le premier définit les opérations \(\lambda\), le second les opérations \(\gamma\).
\[
\begin{tikzcd}[column sep=tiny]
A \arrow[r,"\lambda_u"] \arrow[d,"\lambda_t"'] & \Ghat^u(A)_o \arrow[d,"\Ghat^u(\lambda_t)"]\\
\Ghat^t(A)_o \arrow[r,"\lambda_u"] & \Ghat^u(\Ghat^t(A)_o)_o
\end{tikzcd}
\tag{3.7.1}
\]
\[
\begin{tikzcd}[column sep=tiny]
A \arrow[r,"\gamma_u"] \arrow[d,"\gamma_t"'] & \Ghat^u(A)_\nabla \arrow[d,"\Ghat^u(\gamma_t)"]\\
\Ghat^t(A)_\nabla \arrow[r,"\gamma'_u"] & \Ghat^u(\Ghat^t(A)_\nabla)_\nabla.
\end{tikzcd}
\tag{3.7.2}
\]

3.8. Soit \(K\) un anneau binômial (2.7)~; de (2.6.1) on déduit \(\forall x\in K\) les relations :
\begin{equation}\tag{3.8.1}
  \gamma_t(x)=(1/(1-t))^x
  =\exp(-x\Log(1-t))
\end{equation}

Soit \(a\) un élément de \(K\)~; la multiplication par \(\gamma_t(a)\) pour la loi \(\nabla\) est caractérisée par sa valeur pour \(1+Xt\), soit :
\[
  \gamma_t(a)\nabla(1+Xt)=\rho_t\left(\lambda_t(a)\circ\frac{1+(X+1)t}{1+t}\right)=\rho_t\left(\left(\frac{1+(X+1)t}{1+t}\right)^a\right)
\]
\begin{equation}\tag{3.8.2}
  \gamma_t(a)\nabla(1+Xt)
  =(1+Xt)^a
  =\exp(a\Log(1+Xt))
  =\sum_{n=0}^{\infty}\binom{a}{n}X^nt^n .
\end{equation}
En particulier, on voit que le \(n\)-ième polynôme universel définissant la multiplication par \(\gamma_t(a)\) pour la loi \(\nabla\) est le polynôme \(R_{n,a}\) (2.9).

